We introduce the syntax, the operational semantics, and the type system of
\hyrql{} along with illustrating examples.

\subsection{Syntax}

\begingroup
\renewcommand{\arraystretch}{1.25}
\begin{table}[t]
  \centering
  \[
    \begin{array}{>{$}r<{$} @{\quad} r @{~\Coloneqq~} l}
      (Values) & v & x \ \, \mid\, \zket \gmid \oket \gmid
      c(v_1, \dots, v_n) \gmid \lbd x t \gmid \letrec f x t
      \gmid \sum_{i=1}^n \alpha_i \cdot v_i\\
      (Terms) & t & x\;
      \begin{array}[t]{l}
        \mid\, \zket \gmid \oket \gmid \qcase{t}{t_0}{t_1} \\
        \mid\, c(t_1, \dots, t_n) \gmid \matchdefault t \\
        \mid\, \lbd x t \gmid \letrec f x t \gmid t_1 t_2
        \gmid \sum_{i=1}^n \alpha_i \cdot t_i \gmid \shape t
      \end{array}
    \end{array}
  \]
  \caption{Syntax of the $\hyrql$ language}
  \label{tab:syntax}
\end{table}
\endgroup

The syntax of Terms and Values in $\hyrql$ is provided in
Table~\ref{tab:syntax}. Terms feature variables, noted $x, y, f, \dots$ and
taken from an infinite countable set of variables. We denote by $\seq{e}$ a
(finite and possible empty) sequence of elements $e_1, \dots, e_n$. Qubits are
introduced through basis states $\zket, \oket$, allowing us to write quantum
conditionals $\qcasedefault t$, which corresponds to the superposition of $t_0$
and $t_1$ controlled by qubit $t$. $\hyrql$ also features classical constructors
$c^s$ consisting of a constructor symbol $c$, and a signature $s$ indicating
the type of the constructor. This will be formalized in
Section~\ref{subsec:types} and will be used only when required.
Constructors can be used in constructor application $c(t_1, \dots, t_n)$
and through pattern-matching, sometimes abbreviated as $\matchdefaultcollapsed
t$. Each constructor symbol $c$ comes with a fixed arity and is always fully
applied. We assume the existence of standard inductive constructors for unit
$()$, tensor products $\otimes$, natural numbers $0$ and $S$, and lists $\nil$
and $::$. For convenience, constructors are sometimes used in an infix
notation. For example, tensor products will be written as $x \otimes y$ and
list constructors as $h::t$.

Higher-order is featured via a standard $\lambda$-abstraction $\lbd x t$, and a
construct for general recursion $\letrec f x t$. Term application is denoted by
$t_1t_2$.

Given amplitudes $\alpha_i \in \mathbb C$, the term $\sum_{i=1}^n \alpha_i
\cdot t_i$ represents a \emph{superposition} of terms $t_i$. In the special
case where $n=2$, we just write $\alpha_1 \cdot t_1+\alpha_2 \cdot t_2$. In a
dual manner, the $\textshape$ construct, introduced in~\cite{GLRSV13},
returns the classic structural information on data containing quantum
information. For example, the shape of a qubit list is a copy of the list
structure without the qubits (see Example~\ref{ex:shape-semantics}).


In order to avoid conflicts
between free and bound variables, we will always work up to $\alpha$-conversion
and use Barendregt's convention~\cite[p. 26]{Bar84} which consists in keeping
all bound and free variable names distinct, even when this remains implicit.
We also define a partial order $\subterm$ on terms, where $s \subterm t$ if $s$
is a \emph{strict subterm} of $t$. We denote $\subtermeq$ as the reflexive
closure of $\subterm$. The \emph{size of a term} is written $\size t$ and is
defined standardly in Definition~\ref{def:size}.

As we work with amplitudes and superpositions, Terms and Values have to be
considered with respect to a vector space
structure. Towards that end, we define an equivalence relation $\equiv$ in
Table~\ref{tab:equiv}. Such definition relies on \emph{equivalence contexts},
defined by the following grammar:
\[
  \begin{aligned}
    (\text{Equiv. contexts}) \quad   \equivcontext \Coloneqq \diamond & \gmid
    \qcasedefault \equivcontext \gmid c(\seq{t_1}, \equivcontext, \seq{t_2})
    \gmid t \equivcontext \gmid \equivcontext t\\
    &\gmid \matchdefaultcollapsed{\equivcontext}  \gmid t + \equivcontext \gmid
    \shape \equivcontext
  \end{aligned}
\]
An equivalence context $\equivcontext$ is thus a term with one hole $\diamond$;
let $\equivcontext[t]$ be the term obtained by filling the hole with
$t$ in $\equivcontext$.
The first rules of Table~\ref{tab:equiv} make the definition of the $\sum$
construct unambiguous/sound, i.e., it corresponds exactly to repeated
applications of the $+$ construct. The last rules of Table~\ref{tab:equiv} are
commutation rules highlighting the linearity of quantum terms.

\begingroup
\renewcommand\displaystyle\textstyle
\newcommand\sep{\\[1.25ex]}
\begin{table}[t]
  \centering
  \[
    \begin{array}{c}
      t_1 + t_2 \equiv t_2 + t_1
      \qquad
      t_1 + (t_2 + t_3) \equiv (t_1 + t_2) + t_3
      \qquad
      1 \cdot t \equiv t
      \qquad
      t + 0 \cdot t' \equiv t
      \sep
      \alpha \cdot (\beta \cdot t) \equiv (\alpha \beta) \cdot t
      \qquad
      \alpha \cdot (t_1 + t_2) \equiv \alpha \cdot t_1 + \alpha \cdot t_2
      \qquad
      \alpha \cdot t + \beta \cdot t \equiv (\alpha + \beta) \cdot t
      \\[1.2em]
      \begin{aligned}
        \qcasedefault {(\sum_{j=1}^m \alpha_j \cdot s_j)} &\equiv
        \sum_{j=1}^m \alpha_j \cdot (\qcasedefault{s_j})
        \sep
        c(\seq{t_1}, \sum_{j=1}^m \alpha_j \cdot s_j,
        \seq{t_2}) &\equiv \sum_{j=1}^m \alpha_j \cdot
        c(\seq{t_1}, s_j, \seq{t_2})
        \sep
        \matchdefaultcollapsed{(\sum_{j=1}^m \alpha_j \cdot s_j)}
        &\equiv \sum_{j=1}^m \alpha_j \cdot (\matchdefaultcollapsed{s_j})
        \sep
        t(\sum_{i=1}^m \alpha_j \cdot s_j) &\equiv \sum_{j=1}^m
        \alpha_j \cdot t\,s_j
        \sep
        \equivcontext[t] &\equiv \equivcontext[t'] \qquad \text{ when
        } t \equiv t'
      \end{aligned}
    \end{array}
  \]
  \caption{Equivalence relation $\equiv\ \subseteq\ \text{Terms}
  \times \text{Terms}$}
  \label{tab:equiv}
\end{table}

\endgroup

\begin{example}[Quantum State and Hadamard Gate]
  \label{ex:had-syntax}
  We can define the orthogonal basis states as $\ket \pm \triangleq
  \frac{1}{\sqrt 2} \cdot \zket \pm \frac{1}{\sqrt 2} \cdot \oket$. The
  Hadamard gate can be encoded in \hyrql{}:
  \[
    \mathtt{Had} \triangleq \lbd x \qcase x {\ket +} {\ket -}
  \]
\end{example}

\begin{example}[Quantum Switch]
  \label{ex:qs-syntax}
  The quantum switch~\cite{CAPV13} is a program which, on two unitaries $U, V$
  and a two-qubit state $\ket{\phi}$, applies $UV$ and $VU$, in superposition,
  to the second qubit, controlled by the value of the first qubit. This can be
  implemented in \hyrql{} by $\QS\ U\ V \ket{\phi}$, where $\QS$ is defined as:
  \[
    \QS{} \triangleq\lbd f\lbd g\lbd q \match q { c \otimes t \to
      \qcase{c}{\zket \otimes f(g\,t)}{ \oket \otimes g(f\,t)}
    }
  \]
\end{example}

\begin{example}[Bounded quantum walk]
  \label{ex:bqwalk}
  \hyrql{} allows one to write recursive programs that are classically and
  quantumly controlled. The program $\mathtt{bqwalk}$ produces all the possible
  quantum walk paths of size at most $n$, starting from a quantum state $q$:
  \[
    \begin{aligned}
      \mathtt{repeat} &\triangleq \letrec g n \match n {
        0 \to \nil,
        S(m) \to \zket :: g\,m
      } \\
      \mathtt{bqwalk} &\triangleq \letrec f q \lbd n \qcasesplit q {
        \zket::(\mathtt{repeat}~n)
      } {
        \oket :: 
        \match n {
          0 &\to \nil \\
          S(m) &\to (f\,(\mathtt{Had}\,\oket)\,m)
        }
      }
    \end{aligned}
  \]
\end{example}

The programs from Examples~\ref{ex:qs-syntax} and \ref{ex:bqwalk} are strictly
more than quantum circuits, as one execution of the program does not provide a
full description of the corresponding behaviour. Thus, \hyrql{} is more than a
circuit description language~\cite{GLRSV13}. We also provide the hybrid quantum
protocol BB84~\cite{BB84} and the Quantum Fourier Transform in
Appendix~\ref{app:additional}.

\subsection{Operational Semantics}
\label{subsec:semantics}

\begingroup
\newcommand{\ptlabel}[1]{
  \begin{prooftree}[center=false, right label template =
      \footnotesize\inserttext,
    template = \footnotesize$\inserttext$] #1
\end{prooftree}}
\newcommand{\ptnolabel}[1]{
  \begin{prooftree}[center=false, right label template = ,
    template = \small$\inserttext$] #1
\end{prooftree}}
\renewcommand{\qquad}{\quad\,\,\,}
\renewcommand{\arraystretch}{2.5}
  \begin{table}[t]
    \[
      \begin{array}{c}
        \ptlabel{
          \infer0[(Qcase$_0$)]{\qcase {\zket} {t_0} {t_1 }\reduces t_0}
        }
        \qquad
        \ptlabel{
          \infer0[(Qcase$_1$)]{\qcase {\oket} {t_0} {t_1 }\reduces t_1}
      }
      \\
      \ptlabel{
        \infer0[(Match)]{\matchdefaultcollapsed{c_j(\seq{v_p})
        } \reduces t_j\{ \seq{v_p}/\seq{x_j} \}}
      }
      \qquad
      \ptlabel{
        \infer0[(Lbd)]{(\lbd x t) v_p \reduces t \{v_p/x\}}
      }
      \\
      \ptlabel{
        \infer0[(Rec)]{
          (\letrec f x t)v_p \reduces t \{\letrec f x t/f, v_p/x\}
        }
      }
      \qquad
      \ptlabel{
        \hypo{
          \sum_{i=1}^n \alpha_i \cdot p_i \in \can
          \setminus \valueset
        }
        \hypo{p_i \zoreduces t_i}
        \infer2[(Can)]{\sum_{i=1}^n \alpha_i \cdot p_i \reduces
        \sum_{i=1}^n \alpha_i \cdot t_i}
      }
      \\
      \ptlabel{
        \infer0[(Shape$_0$)]{\shape {\ket 0} \reduces \vunit}
      }
      \qquad
      \ptlabel{
        \infer0[(Shape$_1$)]{\shape {\ket 1} \reduces \vunit}
      }
      \qquad
      \ptlabel{
        \infer0[(Shape$_c$)]{\shape {c^s(\seq{v_p})}
        \reduces   c^{\shape{s}}(\shape{\seq{v_p}})}
      }
      \\
      \ptlabel{
        \hypo{\sum_{i=1}^n \alpha_i \cdot v_i \in \can}
        \infer1[(Shape$_s$)]{\shape{\sum_{i=1}^n \alpha_i \cdot
        v_i} \reduces \shape{v_1}}
      }
      \qquad
      \ptlabel{
        \hypo{t \reduces t'}
        \infer1[(Shape$_\evalcontext$)]{\shape t \reduces \shape{t'}}
      }
      \\
      \ptlabel{
        \hypo{p \reduces t}
        \infer1[($\evalcontext$)]{\evalcontext[p] \reduces \evalcontext[t]}
      }
      \qquad
      \ptlabel{
        \hypo{t \equiv t_1}
        \hypo{t_1 \reduces t_1'}
        \hypo{t_1' \equiv t'}
        \infer3[(Equiv)]{t \reduces t'}
      }
  \end{array}\]
  \caption{Reduction rules of the language}
  \label{tab:reduction}
\end{table}
\endgroup

In this section, we define the call-by-value operational semantics of $\hyrql$.
As we are interested in resource analysis, we want to avoid reducing
non-meaningful terms, e.g., a reduction of $t$ in $s + 0 \cdot t \ \equiv\ s$
is meaningless. Towards that end, we introduce \emph{canonical forms}
(Definition~\ref{def:canonical}), which take care of removing ill-behaved terms
before applying a reduction. It relies on \emph{pure terms}, which are defined
by the following grammar:
\[
  \begin{aligned}
    (\text{Pure terms}) \quad  p  \Coloneqq x &\gmid \zket \gmid \oket
    \gmid \qcasedefault p  \gmid c(p_1, \dots, p_n) \\
    &\gmid \matchdefaultcollapsed p \gmid \lbd x t \gmid \letrec f x t
    \gmid tp \gmid \shape t
  \end{aligned}
\]
A \emph{pure value}, denoted $v_p$, is a pure term that is a value.
\begin{definition}[Canonical form]
  \label{def:canonical}
  A canonical form of a term $t$ is any term $t'$ such that
  $t\equiv t'$ and $t'=\sum_{i=1}^n \alpha_i \cdot p_i $, where $p_i$ are
  pure terms, $\alpha_i \neq 0$ and $p_i \equiv p_j$ implies  $i = j$. The
  set of canonical forms is denoted by $\can$.
\end{definition}

The canonical form always exists and is unique for well-typed terms
(Lemma~\ref{lem:canonical-typed}). Note that the naming \emph{pure terms} comes
from~\cite{DCM22} and is not related to \emph{pure states} in quantum computing,
as pure terms could semantically yield a mixed quantum state.

We also define \emph{evaluation contexts} by the following grammar:
\begin{align*}
  (\text{Evaluation contexts}) \quad \evalcontext \Coloneqq \, \diamond
  &\gmid \qcase{\evalcontext}{t_0}{t_1} \gmid c(\seq{p},
  \evalcontext, \seq{v_p}) \\
  &\gmid \matchdefaultcollapsed{\evalcontext} \gmid t \evalcontext \gmid
  \evalcontext v_p
\end{align*}
Again, let  $\evalcontext[t]$ be the term obtained by filling the
hole $\diamond$ with $t$ in $\evalcontext$.

The operational semantics of $\hyrql$ is described in Table~\ref{tab:reduction}
as a reduction relation $\reduces \ \subseteq
\ \text{Terms} \times \text{Terms}$, where $\{t/x\}$ denotes the
standard substitution of a variable $x$ by a term $t$. The reduction
implements a \emph{call-by-value} strategy.

In the rule (Can), $t\zoreduces t'$  holds if either $t\reduces t'$
or, $\neg(\exists t'',\ t \reduces t'')$ and $t' = t$. Intuitively,
it means that we apply one reduction for each element of a
superposition, when possible.
In the rule (Shape$_c$), $\shape{\seq v}$ is syntactic sugar for $\shape{v_1},
\dots, \shape{v_n}$, and the shape of a signature $\shape s$ is defined
formally in Definition~\ref{def:shape-type}.
In all rules of Table~\ref{tab:reduction}, except for the (Can) and
(Equiv) rules, the left-hand-side term is a pure term (i.e., not a
superposition). This implies that any summation must be
expressed as a canonical form, which avoids reducing
a subterm with a null amplitude, or more generally to reduce two
identical terms that would sum up to $0$.  The goal is to avoid
reductions that have no physical meaning, as this would invalidate
all resource analysis results.

We define $\reduces^*$ as the reflexive and transitive closure of
$\reduces$.
For $k \in \mathbb{N}$, we also write $t \mostreduces k t'$, when $t$
reduces to $t'$ in at most $k$ steps.
A term $t$ \emph{terminates}, if any chain of reduction starting
from $t$ reaches a value, meaning that $t \mostreduces k v$ holds for some $k$.

\begin{example}
  \label{ex:had-semantics}
  Let us consider the $\mathtt{Had}$ program from Example~\ref{ex:had-syntax},
  we can check that it gives the desired result when applying it to $\zket$:
  \[
    \begin{aligned}
      \mathtt{Had} \zket &\reduces \qcase \zket {\ket +}{\ket -}
      &&\text{via (Lbd)} \\
      &\reduces \ket + &&\text{via (Qcase$_0$)}
    \end{aligned}
  \]
  Since the Hadamard gate is its own inverse we can recover $\ket0$.
  Given the following term  $h_v= \qcase{v} {\ket+}{\ket -}$,
  we have the following:
  \[
    \mathtt{Had} \ket + \equiv \frac 1 {\sqrt 2} \cdot \mathtt{Had} \zket +
    \frac 1 {\sqrt 2} \cdot \mathtt{Had} \oket
    \reduces \frac{1}{\sqrt 2} \cdot h_0
    + \frac{1}{\sqrt 2} \cdot h_1 \reduces \frac{1}{\sqrt 2} \cdot \ket + +
    \frac{1}{\sqrt 2} \cdot \ket - \equiv \ket 0
  \]
  Note that the reductions are performed in parallel on each term of the
  superpositions using rule (Can) of Table~\ref{tab:reduction}.
 \end{example}

\begin{example}
  \label{ex:shape-semantics}
  Given a qubit list $l = \zket :: \oket :: \ket + :: \nil$, one can verify that
  $\shape l \sreduces \vunit :: \vunit :: \vunit :: \nil$ (note that $::$ and
  $[\ ]$ use distinct signatures in $l$ and the reduct of $\shape l $).  The
  result is a list of same length as $l$, with no quantum data anymore, hence
  it can be treated non-linearly. This operation is independent of the value of
  each qubit in the list.
\end{example}

It is worth mentioning that while we can write qubit streams in \hyrql{}, the
call-by-value strategy will require it to be evaluated before doing any
operation on such data.

\subsection{Type System}
\label{subsec:types}

\textbf{Types and Contexts. }
Types in \hyrql{} are provided by the following grammar:
\[
  \text{(Types)} \qquad T \Coloneqq \qbit \gmid B \gmid T \linfunc T
  \gmid T \nonlinfunc T
\]
$\qbit$ is the type for qubits and $\bvar$ is a \emph{constructor
type} from a fixed set $\bset$.
A \emph{basic type}, noted $\bkvar$, is either a qubit type, or a
constructor type $\bvar$, i.e., $\bkvar \in \{\qbit\} \cup \bset$.
Each constructor type is defined by its set of constructors:
\[
  \cons B \triangleq \{c^s \gmid s = \signdefault\}
\]

A constructor symbol $c$ comes with a signature $s$ of arity $n$, defining the
inputs and output types of $c$. We may write $c^\bvar$ when $s$ is of arity
$0$. By construction, $\cons B \cap \cons{B'} = \emptyset$ when $B \neq B'$.
In what follows, we will consider the
following constructor types: a unit type $\unit$ with a unique constructor
$\vunit^\unit$; tensor types $\bkvar_1 \otimes \bkvar_2$ for given
types $\bkvar_1, \bkvar_2$ with a unique constructor
$\otimes^{\sign{\bkvar_1, \bkvar_2}{\bkvar_1 \otimes \bkvar_2}}$;
natural numbers $\nat$ with constructors $0^\nat$ and 
$S^{\sign \nat \nat}$; lists $\typelist \bkvar$ for a given type $\bkvar$,
with constructors $\nil^{\typelist \bkvar}$ and $::^{\sign {\bkvar, \typelist
\bkvar} {\typelist \bkvar}}$. We consider that $\bset$ always contains the unit
type $\unit$.

The language also features two distinct \emph{higher-order} types
for linear and non-linear functions, denoted respectively
by $\linfunc$ and $\nonlinfunc$.

In the typing discipline, it will be useful to distinguish between
types based on whether or not they contain quantum data. The latter
must follow the laws of quantum mechanics (no-cloning), whereas
classical types are more permissive.
Towards that end, we define inductively the set of \emph{quantum constructor
types} $\bqset$ by 
\[
  \bqset \triangleq \{\bvar \in \bset \gmid
    \exists\,c^{\signdefault}, \exists\,1 \leq
i \leq n, \bkvar_i \in \bqset \cup \{ \qbit\}\}. 
\]
A quantum
constructor type $\bvar_\qvar \in \bqset$ has (at least) a
constructor symbol with a type $\qbit$ or (inductively) a quantum
constructor type in its signature.
The set $\bcset$ of \emph{classical constructor types} $\bvar_\cvar$
is defined by $\bcset \triangleq \bset \setminus \bqset$.
For example, $\unit,\ \nat \in \bcset$ whereas $\typelist{\qbit} \in \bqset$.
This allows us to split the types between \emph{quantum types} $\qvar$ and
\emph{classical types} $\cvar$:
\[
  \qvar \ \Coloneqq \ \qbit \gmid \bqvar
  \qquad
  \cvar \ \Coloneqq \ \bcvar \gmid T \linfunc T \gmid T \nonlinfunc T
\]
It is now assumed that the inputs of each constructor are ordered, containing
first variables of classical type, then variables of quantum type, enabling a
hybrid behaviour.

\emph{Typing contexts} are defined as follows:
\[
  \Gamma, \Delta \quad \Coloneqq \quad \varnothing \gmid \set{x : T}
  \gmid \set{\shapemarker{x : T}} \gmid \Gamma \cup \Delta,
\]
where $\shapemarker{x : T}$ is called a \emph{boxed variable},
indicating that $x$ is a linear variable captured by a \textshape{} construct.
We define the \emph{domain} of a context $\Delta$ as the set of its variables,
i.e., $\typedom \Delta \triangleq \set{x \gmid \exists\,T,\,x : T
\in \Delta \vee \shapemarker{x : T} \in \Delta}$.
Whenever we write $\Gamma, \Delta$ or
$\Gamma;\Delta$, it is assumed that $\Gamma$ and $\Delta$ are compatible, that
is, $\typedom \Gamma \cap \typedom \Delta = \emptyset$.
We use the shorthand notation $\vartyping {\seq x} {\seq T}$ for $\vartyping
{x_1}{T_1}, \dots, \vartyping {x_n}{T_n}$; and $\shapemarker \Delta$ for
$\shapemarker{x_1 : T_1}, \dots, \shapemarker{x_n : T_n}$.

\medskip
\noindent
\textbf{Typing rules. }
A \emph{typing judgment} is written $\typing \Gamma \Delta t T$, describing
that the term $t$ is \emph{well-typed}, with type $T$, under \emph{non-linear
context} $\Gamma$ and \emph{linear context} $\Delta$.
The typing rules of the language are defined in Table~\ref{tab:typing}.
A well-typed term $\typing \Gamma \Delta t T$ is \emph{closed} if $\Gamma =
\Delta = \emptyset$.
A well-typed closed term $t$ is said to \emph{terminate over any input}, if for
any well-typed closed value $v_1, \dots, v_n$ terminating over any input such
that $tv_1\dots v_n$ is a well-typed closed term of type $\bkvar$, $tv_1 \dots
v_n$ terminates.

\begingroup
\newcommand{\ptlabel}[1]{
  \begin{prooftree}[center=false, right label template = \small\inserttext,
    template = \small$\inserttext$] #1
\end{prooftree}}
\newcommand{\ptnolabel}[1]{
  \begin{prooftree}[center=false, right label template = ,
    template = \small$\inserttext$] #1
\end{prooftree}}
\renewcommand{\arraystretch}{3}
\renewcommand\pt\ptlabel
\begin{table}[t]
  \[
    \begin{array}{c}
      \pt{
        \infer0[(ax)]{\typing \Gamma {\vartyping{x}{T}} x T}
      }
      \qquad
      \pt{
        \infer0[(ax$_c$)]{\typing {\Gamma, \vartyping x C} {} x C}
      }
      \qquad
      \pt{
        \infer0[(ax$_0$)]{\typing \Gamma {} {\zket} \qbit}
      }
      \\
      \pt{
        \infer0[(ax$_1$)]{\typing \Gamma {} {\oket} \qbit}
      }
      \quad
      \ptlabel{
        \hypo{\typing \Gamma \Delta t \qbit}
        \hypo{\typing \Gamma {\Delta'} {t_0} Q}
        \hypo{\typing \Gamma {\Delta'} {t_1} Q}
        \hypo{t_0 \perp t_1}
        \infer4[(qcase)]{\typing \Gamma {\Delta, \Delta'} {\qcase
        t {t_0} {t_1}} Q}
      }
      \\
      \pt{
        \hypo{s = \signdefault}
        \hypo{\typing \Gamma {\Delta_i} {t_i} {\bkvar_i}}
        \infer2[(cons)]{\typing \Gamma {\Delta_1, \dots, \Delta_n}
        {c^s(t_1, \dots, t_n)} \bvar}
      }
      \\
      \ptlabel{
        \hypo{\cons \bvar = \set{c_i^{s_i}}_{i=1}^n}
        \hypo{s_i = \sign{\seq{C_i}, \seq{Q_i}} \bvar}
        \hypo{\typing {\Gamma_1}{\Delta_1} t \bvar}
        \hypo{\typing
          {\Gamma_2, \vartyping{\seq{y_i}}{\seq{C_i}}}
          {\Delta_2, \vartyping{\seq{z_i}}{\seq{Q_i}}}
          {t_i} \bkvar
        }
        \infer4[(match)]{
          \typing {\Gamma_1, \Gamma_2} {\Delta_1, \Delta_2} {
            \textmatch_{1 \leq i \leq n}\,t\,\{
              c_i^{s_i}(\seq{y_i}, \seq{z_i}) \to t_i
            \}
          } \bkvar
        }
      }
      \\
      \pt{
        \hypo{\typing \Gamma {\Delta, \vartyping x T} t T'}
        \infer1[(abs)]{\typing \Gamma \Delta {\lbd x t} {T \linfunc T}'}
      }
      \qquad
      \pt{
        \hypo{\typing {\Gamma, \vartyping x C} \Delta t T}
        \infer1[(abs$_c$)]{\typing \Gamma \Delta {\lbd x t} {C \nonlinfunc T}}
      }
      \qquad
      \pt{
        \hypo{\typing {\Gamma, \vartyping f T} {} {\lbd x t} T}
        \infer1[(rec)]{\typing \Gamma {} {\letrec f x t} T}
      }
      \\
      \qquad
      \pt{
        \hypo{\typing \Gamma \Delta {t_1} {T \linfunc T'}}
        \hypo{\typing \Gamma {\Delta'} {t_2} T}
        \infer2[(app)]{\typing \Gamma {\Delta, \Delta'} {t_1t_2} {T'}}
      }
      \qquad
      \pt{
        \hypo{\typing \Gamma \Delta {t_1} {C \nonlinfunc T}}
        \hypo{\typing \Gamma {} {t_2} C}
        \infer2[(app$_c$)]{\typing \Gamma \Delta {t_1t_2} T}
      }
      \\
      \hypertarget{typesup}{
      \ptlabel{
        \hypo{\typing \Gamma \Delta {t_i} \qvar}
        \hypo{\sum_{i=1}^n \size{\alpha_i}^2 = 1}
        \hypo{\forall i \neq j,\ t_i \perp t_j}
        \infer3[(sup)]{\typing \Gamma \Delta{ \sum_{i=1}^n \alpha_i
        \cdot t_i} \qvar}
      }
      }
      \quad
      \ptlabel{
        \hypo{\typing \Gamma \Delta t \bkvar}
        \infer1[(shape)]{\typing {\Gamma, \shapemarker \Delta} {}
        {\shape{t}} {\shape \bkvar}}
      }
      \\
      \ptlabel{
        \hypo{\typing {\Gamma,\shapemarker{\vartyping x \bkvar}}
        {\Delta, \vartyping y \bkvar} t T}
        \infer1[(contr)]{\typing \Gamma {\Delta, \vartyping y \bkvar}
        {t\{y/x\}} T}
      }
      \qquad
      \ptlabel{
        \hypo{\typing \Gamma \Delta t T}
        \hypo{t \equiv t'}
        \infer2[(equiv)]{\typing \Gamma \Delta {t'} T}
      }
  \end{array}\]
  \caption{Typing rules of the language}
  \label{tab:typing}
\end{table}
\endgroup

Typing rules rely on an external \emph{orthogonality} predicate on terms. This
uses the notion of $\Delta$-\emph{context substitutions}, which, for a given
$\Delta$, denote the substitutions $\sigma$ such that for any $x : T$ with
$\vartyping x T \in \Delta$ or $\shapemarker{\vartyping x T} \in \Delta$,
$ \sigma(x)=v$, where $v$ is a well-typed closed value of type $T$
that terminates over any given input. It also requires the definition of an
an \emph{inner product} on values with a canonical form, defined
below, where $\kron x y$ is the Kronecker symbol.
\[
  \langle v, w \rangle \triangleq \sum_{i=1}^n \sum_{j=1}^m \alpha_i \beta_j^*
  \kron{v_i}{w_j}, \quad \text{ where }
  v \equiv \sum_{i=1}^n \alpha_i \cdot v_i \in \can \text{ and }
  w \equiv \sum_{j=1}^m \beta_j \cdot w_j \in \can
\]
\begin{definition}[Orthogonality]
  \label{def:orthogonality}
  Let $\typing \Gamma \Delta t Q$ and $\typing \Gamma \Delta {t'} Q$ be
  two well-typed terms. $t$ and $t'$ are \emph{orthogonal}, written $t \perp
  t'$, if for any $\Gamma \cup \Delta$-context substitution $\sigma$:
  \begin{itemize}
    \item $ \sigma(\shape t) \sreduces v$, $\sigma(\shape{t'})
      \sreduces v'$ and $v = v'$;
    \item $\sigma(t) \sreduces v \in \can$,
      $\sigma(t') \sreduces v' \in \can$,
      and $\langle v, v' \rangle = 0$.
  \end{itemize}
\end{definition}

These conditions imply that the reduced values share the same classical
structure, and that their quantum data is orthogonal. For example, orthogonal
lists must be of the same size. This is used in the typing rules
(qcase) and (sup) of Table~\ref{tab:typing} to ensure the feasibility of typed terms, similarly to
existing languages \cite{DCM22,Lem24,DLPZ25}. Note that the rules requiring
this predicate apply it only to well-typed terms, thus the type system is not
cyclic. While an alternative condition could have been given with dependent
types, this gives a uniform process that is defined for any type, including
types that could be defined later by a programmer, and a simpler type system.
As we will prove that values have a unique canonical form
(Lemma~\ref{lem:canonical-typed}), the inner-product condition is sound.

We also define the \emph{shape of a basic type} below.
\begin{definition}
  \label{def:shape-type}
  The shape of a basic type is defined inductively as $\shape \qbit \triangleq
  \unit$, and $\shape \bvar $ is the type defined by the constructor set  
  \[
    \cons{\shape \bvar} \triangleq \set{
      c^{\shape{s}}
      \gmid c^{s} \in \cons \bvar
    }
  \]
  where the shape of a signature $s=\bkvar_1,\ldots,\bkvar_n \to \bvar$ is defined as 
  \[
  \shape{s} \triangleq \shape{\bkvar_1},\ldots,\shape{\bkvar_n} \to \shape{\bvar}.
  \]
\end{definition}
In particular, the shape of a list of qubits is a list of units, i.e.,
$\shape{\typelist{\qbit}} = \typelist{\unit}$, which is coherent with the
reduction rules in Table~\ref{tab:reduction}. By construction, it always holds
that $\shape \bkvar \in \bcset$. This is used in the typing rule (shape) of
Table~\ref{tab:typing}, which extracts the classical structure of a quantum
term, boxes any variable in $\Delta$, and puts it in the non-linear context.
The judgment obtained has an empty linear context and thus can be used
classically (i.e., duplicated or discarded). A boxed variable $x$ can be
removed using rule (contr) of Table~\ref{tab:typing}, by associating with a
linear variable of same type $y$, allowing to read both the quantum data and
the classic structural information of $y$ (see Example~\ref{ex:len-typing}).


\medskip
\noindent
\textbf{Typing derivations. }
We write $\pi \triangleright \typing \Gamma \Delta t T$ for the tree of root
$\typing \Gamma \Delta t T$ derived using the rules of Table~\ref{tab:typing}.
We also denote by $\pi \rightslice \mathrm{(r)}[t_1, \dots, t_n]$ the subtree
of $\pi$ rooted by (r) and whose premises are $t_1, \dots, t_n$.

\begin{example}\label{ex:len-typing}
  The length of a list can be computed by the following term:
  \[
    \mathtt{len} \triangleq \letrec fx \match x {
      \nil \to 0, h::t \to S(f\,t)
    }
  \]
  Notice that we discard $h$, thus the list is not used linearly. Therefore, we
  can only type this as a non-linear function, i.e., $\closedtyping
  {\mathtt{len}} \typelist B \nonlinfunc \nat$ for a given $B \in \bcset$.
  However, we can still compute the length of a qubit list with $\lbd x (x
  \otimes \mathtt{len}(\shape x))$, whose typing derivation $\pi$ is written
  below. For conciseness, we hide the signature of $\otimes$.
  \[
    \begin{prooftree}[
        center=false, right label template = \footnotesize\inserttext, rule
        margin = 1.2ex, template = \footnotesize$\inserttext$
      ]
      \infer0[(ax)]{
        \typing{\shapemarker{y : \typelist \qbit}}{x : \typelist
        \qbit}{x}{\typelist \qbit}
      }
      \infer0[(ax)]{\typing{}{y : \typelist \qbit}{y}{\typelist \qbit}}
      \infer1[(shape)]{
        \typing{\shapemarker{y : \typelist \qbit}}{}{\color{red}\shape y}{\typelist \unit}
      }
      \infer1[\color{red}(app$_c$)]{
        \typing{\shapemarker{y : \typelist
        \qbit}}{}{\mathtt{len}(\shape y)}{\nat}
      }
      \infer2[(cons)]{
        \typing{\shapemarker{y : \typelist \qbit}}{x : \typelist \qbit}{x
        \otimes \mathtt{len}(\shape y)}{\typelist \qbit \otimes \nat}
      }
      \infer1[(contr)]{
        \typing{}{x : \typelist \qbit}{x \otimes \mathtt{len}(\shape
        x)}{\typelist \qbit \otimes \nat}
      }
      \infer1[(app)]{
        \typing{}{}{\lbd x (x \otimes \mathtt{len}(\shape
        x))}{\typelist \qbit \linfunc \typelist \qbit \otimes \nat}
      }
    \end{prooftree}
  \]
  As $\shape y$ is of classical constructor type, it can be used as an input of
  $\mathtt{len}$, allowing us to compute the length of $y$, given by the
  subtree $\pi \rightslice {\color{red}(\mathrm{app}_c) [\shape y]}$. The rule
  (contr) removes the marker introduced by (shape), by substituting $y$ with a
  variable of same type, $x$, thus yielding a closed term. Note that while $x$
  appears twice in the final term, the quantum information of $x$ has a linear
  behaviour, and only the classical structure is used non-linearly.
\end{example}
